S6 receives original RGB8 images and a double-precision homography
$H_{B\rightarrow A}$. The caller owns pitched buffers, workspace and stream.
Geometry preparation rejects singular or non-finite transforms, a projective
horizon crossing B, and excessive canvas bounds. Integer pixel centres define
an inclusive canvas enclosing A and projected B. At canvas coordinate $(x,y)$,
the renderer samples $p_A=(x+x_0,y+y_0,1)^{\mathsf T}$ and
$p_B\sim H^{-1}p_A$, using FP64 coordinates and FP32 bilinear colour arithmetic.
Validity is independent of RGB intensity. Source-only pixels retain their
warped colour, while pixels outside both images have zero RGB and union mask.

For source $i$, the feather weight is its normalized distance to the closest
image edge,
\begin{equation}
d_i=\min\!\left\{\frac{u_i+1/2}{W_i},\frac{W_i-1/2-u_i}{W_i},
                 \frac{v_i+1/2}{H_i},\frac{H_i-1/2-v_i}{H_i}\right\}.
\end{equation}
Overlap colour is $(d_A C_A+d_B C_B)/(d_A+d_B)$. The CUDA feather path
fuses inverse warping and blending without temporary workspace.

Multi-band blending follows Gaussian/Laplacian pyramids~\cite{burt1983spline}.
The source with larger $d_i$ owns an overlap pixel, with ties assigned to A.
A pull--push extension fills invalid colour samples for filtering while
preserving the original validity mask. Reduction filters with
$[1,4,6,4,1]/16$ and subsamples at even coordinates. Expansion inserts zeros,
filters with gain four and crops to the parent shape. Borders use reflect-101.
Laplacian bands from both sources are combined using Gaussian ownership
weights, then reconstructed from coarse to fine. Requested depths three and
five count reductions, yielding four and six stored levels where dimensions
permit. The final mask restores source-only and uncovered pixels. Signed bands
retain their range until final RGB8 quantization. Kernels evaluate the filter
directly. Allocation and synchronization are managed by the caller.
